\documentclass [
%	draft=true, 					%% Entwurfsstadium
%	final, 							%% fertiges Dokument
	fontsize=12pt, 					%% Schriftgroesse
	a4paper, 						%% Papiergroesse
%	titlepage=true, 				%% Titelseite auf extra Seite
	oneside, openany, 				%% 
%	twoside, openright, 			%%
%%------------------------------------------------------------------------------
%% Parameter fuer Verzeichnisse:
%%------------------------------------------------------------------------------ 
	toc=bibliography, 				%% Bibliographie ins TOC
%	toc=bibliographynumbered, 		%% Bibliographie im TOC nummeriert
%	toc=chapterentrywithdots, 		%% Kapiteleintrag, Seitenzahl mit Punktlinie
%	toc=chapterentrywithoutdots, 	%% Kapiteleintrag, Seitenzahl ohne Punktlinie
%	toc=flat, 						%% TOC in tabelarischer Form
%	toc=graduated, 					%% TOC in hierarchischer Form
%	toc=numberline, 				%% buendig mit nummeriertem Eintrag
%	toc=nonumberline, 				%% buendig mit Nummer
	toc=index, 						%% Index ins TOC (idxtotoc)
	toc=listof, 					%% Abbildungs- und das Tabellenverzeichnis
%	toc=listofnumbered, 			%% Abbildungs- und das Tabellenverzeichnis mit Nummer
%	toc=sectionentrywithdots, 		%% Abschnittseintrag mit Punktlinie
%	toc=sectionentrywithoutdots, 	%% Abschnittseintrag ohne Punktlinie
%%------------------------------------------------------------------------------
%	chapterprefix=true, 			%% Default: false
%	appendixprefix=true, 			%% Default: false
] {scrreprt} 					%% Dokumentklasse [scrartcl|scrreprt|scrbook]
%%==============================================================================
%% Document metadata
%%==============================================================================
%% general metadata:
\newcommand{\myname}{Max Mustermann}
\newcommand{\mymail}{\href{mailto:Max.Mustermann@example.com}{Max.Mustermann@example.com}}
\newcommand{\myauthor}{\myname \\\mymail} 
\newcommand{\mytitle}{TITLE}  		%% also used for PDF metadata (hyperref)
\newcommand{\mysubject}{SUBJECT} 	%% also used for PDF metadata (hyperref)
\newcommand{\mykeywords}{KEYWORDS} 	%% also used for PDF metadata (hyperref)
\newcommand{\mypdfauthor}{\myname}	%% used for PDF metadata (hyperref)
%
\newcommand{\myerstpruefer}{Prof. Dr. Pruefer} 			%% used in Titlepage
\newcommand{\myzweitpruefer}{Prof. Dr. Zweitpruefer} 	%% used in Titlepage
\newcommand{\myabgabedatum}{\today} 					%% used in Titlepage and PDF metadata (hyperref)
%
% ** Packages ***************************************
\usepackage {fixltx2e} %korrigiert einige Probleme in LaTeX
\usepackage {fix-cm} %Behebung von Fehlern in den Standard-Schriften
%\usepackage [ansinew] {inputenc} %Benutzung von Umlauten und Sonderzeichen im Quelltext (Windows-Codierung)
\usepackage [utf8]{inputenc} %Benutzung von Umlauten und Sonderzeichen im Quelltext (Mac-Codierung)
\usepackage [T1] {fontenc} %Zeichenbelegung der verwendeten Ausgabeschrift
\usepackage {textcomp} %Sonderzeichen
%
\usepackage [english, ngerman] {babel} %Sprachuntersttzung
%
\usepackage {graphicx} 	% Graphik-Paket
\usepackage {caption} 	% customizing captions of floating environments
%
\usepackage {fixmath} %internationaler Standard fr mathematische Formeln
\usepackage {amsmath} %erweiterter Mathematik-Modus
\usepackage {amssymb} %mathematische Symbole
%
\usepackage[
	draft=false,
	automark, 			% 
	autooneside, 		% fuer einseitigen Druck
	footsepline, 		% Linie oberhalb Fusszeile
	headsepline, 		% Linie unterhalb Kopfzeile
	plainheadsepline, 	% Linie unterhalb Kopfzeile auf Kapitelanfangsseiten
	plainfootsepline, 	% Linie oberhalb Fusszeile auf Kapitelanfangsseiten
	]{scrlayer-scrpage} % Formatierung der Kopf- und Fuzeilen
%
\usepackage{geometry} % Seitenlayout veraendern
%
\usepackage[
	draft=false,
]{hyperref} %pdf-Optionen
%
\usepackage[
	draft=false,
]{bookmark}
%
\usepackage{verbatim} %Verbatim-Umgebung nutzen
\usepackage{listings} %Package fuer Quellcodelistings
\usepackage[ %
%	toc, 		% puts a header into ToC
%	page, 		% puts a title into document
%	titletoc, 	% adds 'Appendix' prefix to titles in ToC
%	header, 	% adds 'Appendis' prefix to pageheaders
]{appendix} 	% Package for \begin{appendices}
%
\usepackage{url} 		% for some natbib styles
\usepackage[ %
	round, 				% (default) for round parentheses
%	square, 			% for square brackets
%	curly, 				% for curly braces
%	angle, 				% for angle brackets
	semicolon, 			% (default) to separate multiple citations with semi-colons
%	comma, 				% to use commas as separators
	authoryear, 		% (default) for author-year citations
%	numbers, 			% for numerical citations
%	super,				% for superscripted numerical citation, as in Nature
%	sort, 				% orders multiple citations into the sequence in which they appear in the list of references
%	sort&compress, 		% as sort but in addition multiple numerical citations are compressed if possible (as 3-6, 15)
%	compress, 			% to compress without sorting
%	sectionbib, 		% redefines \thebibliography to issue \section* instead of \chapter*; 
%	nonamebreak, 		% keeps all author's names in a citation on one line
]{natbib}				% Natural Sciences Citations and References
%
\usepackage[ %
	autostyle=true, 	% default: tryonce
	german=quotes, 		% syle variants for german: quotes, guillements, swiss
	english=british, 	% syle variants for english: american, british
	french=quotes, 		% syle variants for french: quotes, quotes*, guillements, guillements*
%	maxlevel=2, 		% default: 2
%	threshold=3, 		% default: 3
	]{csquotes}
%
\usepackage{relsize} 	% requiered by acronym option smaller
\usepackage[
	printonlyused, 		%
%	footnote, 			% makes the full name apear as footnote
%	withpage, 			% marker which tells to print page numbers
%	smaller				% lets acronyms appear a bit smaller than surrounding text
	]{acronym} 			% 
%
\usepackage{makeidx} 	% 
%
\usepackage{array} % for Titlepage 	% Pakete einbinden
%*******************************************************
% Layout
%*******************************************************
%\automark[section]{chapter} 	% chapter-klasses default
%\automark[subsection]{section}	% others default
%
\clearscrheadfoot% scrheadings und scrplain komplett leeren
%% Kopfzeile
%% --- \leftmark \rightmark \headmark \pagemark ---
%\lehead []{} 			% left aligned even headings
%\cehead []{} 			% centered even headings
%\rehead []{} 			% right aligned even headings
%\lohead []{} 			% left aligned odd headings
%\cohead []{} 			% centered odd headings
%\rohead []{} 			% right aligned odd headings
\ihead []{} 				% inside headings
\chead []{} 				% centered headings
\ohead []{\headmark} 		% outside headings
%% Fusszeile
%\lefoot []{} 			% left aligned even headings
%\cefoot []{} 			% centered even headings
%\refoot []{} 			% right aligned even headings
%\lofoot []{} 			% left aligned odd headings
%\cofoot []{} 			% centered odd headings
%\rofoot []{} 			% right aligned odd headings
\ifoot [\mytitle]{\mytitle} 		% inside headings
\cfoot []{} 				% centered headings
\ofoot [\pagemark]{\pagemark} 		% outside headings
\pagestyle {scrheadings} %Stil der Kopf- und Fuzeilen
%
\addtokomafont {captionlabel} {\sffamily\bfseries} %Unterschriftentitel fett und serifenlos
\setcapindent {1em} %hngender Einzug ab zweiter Zeile
%
\renewcommand {\thefigure} {\thesection.\arabic{figure}} %Sektionsnummer mit in Abbildungsnummer
\numberwithin {figure} {section} %Abbildungsnummer mit jeder Sektion zurcksetzen
%
\geometry{a4paper, top=27mm, left=30mm, right=20mm, bottom=35mm, headsep=10mm, footskip=12mm} 		% Layout
\makeatletter
\@ifclassloaded{scrreprt}{
\newcommand\frontmatter{%
    \cleardoublepage
  %\@mainmatterfalse
  \pagenumbering{Roman}
  \renewcommand{\thechapter}{\Roman{chapter}}
  \renewcommand{\thesection}{\Roman{section}} 
  \renewcommand{\thesubsection}{\thesection.\Roman{subsection}}}
%
\newcommand\mainmatter{%
    \cleardoublepage
 % \@mainmattertrue
  \pagenumbering{arabic}
  \renewcommand{\thechapter}{\arabic{chapter}}
  \renewcommand{\thesection}{\arabic{section}} 
  \renewcommand{\thesubsection}{\thesection.\arabic{subsection}}}
%
\newcommand\backmatter{%
    \cleardoublepage
 % \@mainmatterfalse
   }
}{}
\makeatother
%
\makeindex
%
\SetCiteCommand{\citep}
 	% Befehle einbinden
\hypersetup {
	plainpages = {false}, %verwendet nur arabische Seiten-Label
	breaklinks = {true}, %Erlaubt Zeilenumbrche in Links
	colorlinks = {false}, %Benutze farbige Links
	pdftitle = {\mytitle}, %Titel der Arbeit
	pdfsubject = {\mysubject}, %Thema der Arbeit
	pdfauthor = {\mypdfauthor}, %Autor der Arbeit
	pdfkeywords = {\mykeywords}, %Stichwrter zur Arbeit
	pdfproducer = {pdfLatex}, %Erzeugt durch
	pdfcreator = {LaTeX (MiKTeX mit TeXnicCenter)} %Erstellt mit
}
%
%\pdfinfo{
%	/Title ()
%	/Author ()
%	/Creator ()
%	/Producer ()
%	/CreationDate (201501011200)
%	/ModDate (201501011200)
%	/Subject ()
%	/Keywords ()
%} 	% hyperref pdf-Optionen
%%=Titelseite=================================\maketitle
\titlehead {} 			%% Kopfzeile a. d. Titelseite
\title {\mytitle} 		%% Title a. d. Titelseite
\subtitle {} 			%% Untertitel a. d. Titelseite
\subject {
	\includegraphics[width=0.6\linewidth]{static/hs_os_iui}\\
	{\large Fakultät Ingenieurwissenschaften und Informatik}\\
	\vspace{2ex}
	\mysubject
} 	%% zusätzlicher Titel a. d. Titelseite
\author {\myauthor} 	%%
\date {} 			%%
\publishers {
	\normalsize
	\begin{tabular}{ll}
	  1. Prüfer &: \myerstpruefer\\
	  2. Prüfer &: \myzweitpruefer\\
	   & \\
	  Abgabedatum &: \myabgabedatum\\
	\end{tabular}
} 			%% Herausgeber/Verlag a. d. Titelseite
%\dedication {} 			%% Widmung nach der Titelseite
%\extratitle {} 			%% Schmutztitel vor der Titelseite
%\uppertitleback {} 		%% Kopzeile nach der Titelseite
%\lowertitleback {} 		%% Fußzeile nach der Titelseite
%%==========================================\Titelseite= 	% options for maketitle
%*******************************************************
% Listings
%*******************************************************
% Use with 
% 	\begin{lstlisting}[style=example, caption={Descriptive Caption Text}, label=DescriptiveLabel]
% 	\end{lstlisting}
% OR
% 	\lstinputlisting[style=example, caption={Descriptive Caption Text}, label=DescriptiveLabel]{hello.c}
%
\lstdefinestyle{example}{ 		%% \lstset{style=example}
  basicstyle=\footnotesize, 
  captionpos=b, 
  breaklines=true, 
  showstringspaces=false, 
  tabsize=2, frame=lines, 
  numbers=left, 
  numberstyle=\tiny, 
  xleftmargin=2em, 
  framexleftmargin=2em
}
%----------------------------------
% Einige Parameter
%----------------------------------
%\lstset{ %
%  backgroundcolor=\color{white},   % choose the background color; you must add \usepackage{color} or \usepackage{xcolor}
%  basicstyle=\footnotesize,        % the size of the fonts that are used for the code
%  breakatwhitespace=false,         % sets if automatic breaks should only happen at whitespace
%  breaklines=true,                 % sets automatic line breaking
%  captionpos=b,                    % sets the caption-position to bottom
%  commentstyle=\color{mygreen},    % comment style
%  deletekeywords={...},            % if you want to delete keywords from the given language
%  escapeinside={\%*}{*)},          % if you want to add LaTeX within your code
%  extendedchars=true,              % lets you use non-ASCII characters; for 8-bits encodings only, does not work with UTF-8
%  frame=single,                    % adds a frame around the code
%  keepspaces=true,                 % keeps spaces in text, useful for keeping indentation of code (possibly needs columns=flexible)
%  keywordstyle=\color{blue},       % keyword style
%  language=Octave,                 % the language of the code
%  otherkeywords={*,...},            % if you want to add more keywords to the set
%  numbers=left,                    % where to put the line-numbers; possible values are (none, left, right)
%  numbersep=5pt,                   % how far the line-numbers are from the code
%  numberstyle=\tiny\color{mygray}, % the style that is used for the line-numbers
%  rulecolor=\color{black},         % if not set, the frame-color may be changed on line-breaks within not-black text (e.g. comments (green here))
%  showspaces=false,                % show spaces everywhere adding particular underscores; it overrides 'showstringspaces'
%  showstringspaces=false,          % underline spaces within strings only
%  showtabs=false,                  % show tabs within strings adding particular underscores
%  stepnumber=2,                    % the step between two line-numbers. If it's 1, each line will be numbered
%  stringstyle=\color{mymauve},     % string literal style
%  tabsize=2,                       % sets default tabsize to 2 spaces
%  title=\lstname                   % show the filename of files included with \lstinputlisting; also try caption instead of title
%} 	% Listings
%
%%==============================================================================
%% Document content 
%%==============================================================================
\begin {document} % hier beginnt das eigentliche Dokument
\selectlanguage{ngerman}
%*****************************************
% Frontmatter
%*****************************************
\frontmatter 								%% scrbook only
%\maketitle
\selectlanguage{ngerman}
\begin{titlepage}\noindent
  \begin{center}
%    \includegraphics[width=0.8\linewidth]{static/hs_os}\\
  	\includegraphics[width=0.8\linewidth]{static/hs_os_iui}\\
  	\vspace*{2cm}
  	\Large
  	\textbf{Fakultät}\\
  	\textbf{Ingenieurwissenschaften und Informatik}\\
  	\vspace*{2cm}
  	\Huge
  	\textbf{\mysubject}\\
  	\vspace*{0.5cm}
  	\large
  	über das Thema\\
  	\vspace*{1cm}
  	\textbf{\mytitle}\\
  	\vspace*{2cm}	
  	\vfill
  	\normalsize
  	\newcolumntype{x}[1]{>{\raggedleft\arraybackslash\hspace{0pt}}p{#1}}
  	\begin{tabular}{x{6cm}p{7.5cm}}
  		\rule{0mm}{5ex}\textbf{Autor:} & \myname\newline \mymail \\ 
  		\rule{0mm}{5ex}\textbf{1. Prüfer:} & \myerstpruefer \\ 
  		\rule{0mm}{5ex}\textbf{2. Prüfer:} & \myzweitpruefer \\ 
  		\rule{0mm}{5ex}\textbf{Abgabedatum:} & \myabgabedatum \\ 
  	\end{tabular} 
  \end{center}
\end{titlepage}
 			%% alternative zu \maketitle
\begin{abstract} 							%% [scrartcl|scrreprt] only
% ----------------------------------------------------------------------------------------------------------
% Abstract
% ----------------------------------------------------------------------------------------------------------
\selectlanguage{ngerman}
\section*{Kurzfassung}
\addcontentsline{toc}{chapter}{Kurzfassung}
%\blindtext

\selectlanguage{english}
\section*{Abstract}
%\blindtext.

\end{abstract} 								%% [scrartcl|scrreprt] only
\selectlanguage{ngerman}
\tableofcontents 		% Inhaltsverzeichnis
\listoffigures 			% Abbildungsverzeichnis
\listoftables 			% Tabellenverzeichnis
\lstlistoflistings 		% Quellcodelistingverzeichnis

\clearpage
\selectlanguage{ngerman}
% ------------------------------------------------------------------------------
% Abkürzungen
% 
% Use: 
% 	\ac{<acronym>} 		first use, full name is printed
% 	\acf{<acronym>}		full name of acronym is printed
% 	\acs{<acronym>}		short name of acronym is printed
% 	\acl{<acronym>} 	expanded acronym without mentioning the acronym
% 	\acp{<acronym>} 	same as \ac, but plural
% 	\acfp{<acronym>} 	same as \acf, but plural
% 	\acsp{<acronym>} 	same as \acs, but plural
% 	\aclp{<acronym>} 	same as \acl, but plural
% 	\iac{<acronym>} 	prefixes \ac with an appropriate indefinite article
% 	\Iac{<acronym>} 	
% 	\...*{<acronym>}	will not be marked as used
% ------------------------------------------------------------------------------
\chapter*{Abk\"urzungsverzeichnis}
\addcontentsline{toc}{chapter}{Abk\"urzungsverzeichnis} 
\begin{acronym}[H2O] % längste Abkürzung steht in eckigen Klammern
	\setlength{\itemsep}{-\parsep} % geringerer Zeilenabstand
%%--------------------------------------------------------------------
% 	\acro{acronym}[short name]{long name}
% 	\acro{acronym}[short name]{\foreignlanguage{english}{}}
%	\acroplural{acronym}[short plural]{long plural}
% 	\acro{acronym}[short name]{full name \acroextra{additional info}}
%%--------------------------------------------------------------------
%	\acro{H2O}[$\mathrm{H_2O}$]{Dihydrogen Monoxide\acroextra{ (Water)}} 
\end{acronym}
%
%*****************************************
% Mainmatter
%*****************************************
\mainmatter 								%% scrbook only
\renewcommand{\thesection}{\thechapter .\arabic{section}} % Nummeriert Kapitel.Section.Subsection
%\cleardoublepage\include{Sections/Section01}
%\cleardoublepage\include{Sections/Section02}
%\cleardoublepage\include{Sections/Section03}
%
%*****************************************
% Backmatter
%*****************************************
\backmatter 								%% scrbook only
\bookmarksetup{startatroot}
\begin{appendices}
%\include{Sections/Anhang}

\end{appendices}
%----Literaturverzeichnis-------------%
\bibliographystyle{dinat} % Literaturverzeichnis nach DIN 1505; [alphadin|natdin|dinat|dcu]
\bibliography{bib}
\pagebreak
%----Index----------------------------%
\printindex
%
\cleardoublepage\thispagestyle{empty}
\begin{center}
	\vspace*{5em}
	\huge\textbf{Erklärung}\\
\end{center}
\vspace{2em}
Hiermit versichere ich, dass ich meine Abschlussarbeit selbständig verfasst und keine anderen als die angegebenen Quellen und Hilfsmittel benutzt habe.

\vspace{4em}
\begin{minipage}{\linewidth}
	\begin{tabular}{p{15em}p{15em}}
		Datum: &  .......................................................\\
		& \centering (Unterschrift)\\
	\end{tabular}
\end{minipage}
%
\end{document} %hier endet das eigentliche Dokument